% questions.tex —— 1990 年上海卷数学 · 题目内容片段（无答案）
% 由 tools/generate_projects.py 依来源总账自 1990/1990_数学真题_上海_试题.md 生成，勿手改（改源 MD 后重跑）


\begin{gktip}
\textbf{转录边界}：依据固定来源单页 A3 试题原件转录。第 1—10 题为填空题，第 11—20 题为选择题，第 21—27 题为解答题。已完成冷启动盲转录核对；不从其他试卷或答案推补内容。

\textbf{待核处}：第 18 题第三个量印作 \texttt{b}，而选项使用 \texttt{c}；第 23 题原文在“点 P”后直接接“直线 $x=2$”，未见“在”字。均按原件照录，不擅自改正。第 15 题已按“角属于第二象限”的终边定义解析，仍仅为独立推导，非来源印载答案。

\textbf{图形说明}：第 26 题原卷含两平面及点线关系图，现附按点线拓扑重绘的 TikZ 矢量示意图；图不按比例，不提交扫描页图片。
\end{gktip}

\gkpart{试题}

\gkpart{一、填空题}

\begin{question}
\gkti{1}函数 $y=\dfrac{\sqrt{x+4}}{x+2}$ 的定义域是\gkblank{}。

\end{question}
\begin{question}
\gkti{2}函数 $y=\arcsin x$（$x\in[-1,1]$）的反函数是\gkblank{}。

\end{question}
\begin{question}
\gkti{3}过点 $(1,2)$ 且与直线 $2x+y-1=0$ 平行的直线方程是\gkblank{}。

\end{question}
\begin{question}
\gkti{4}已知圆柱的轴截面是正方形，它的面积是 $4\,\mathrm{cm}^2$，那么这个圆柱的体积是\gkblank{}$\,\mathrm{cm}^3$。（结果中保留 $\pi$）

\end{question}
\begin{question}
\gkti{5}在 $\triangle ABC$ 中，已知 $\cos A=-\dfrac35$，则 $\sin\dfrac A2=$\gkblank{}。

\end{question}
\begin{question}
\gkti{6}设复数 $\omega=\cos\dfrac{2\pi}{5}+i\sin\dfrac{2\pi}{5}$，则 $\omega+\omega^2+\omega^3+\omega^4+\omega^5$ 的值是\gkblank{}。

\end{question}
\begin{question}
\gkti{7}已知圆锥的中截面周长为 $a$，母线长为 $l$，则它的侧面积等于\gkblank{}。

\end{question}
\begin{question}
\gkti{8}已知 $(x+a)^7$ 的展开式中，$x^4$ 的系数是 $-280$，则实数 $a=$\gkblank{}。

\end{question}
\begin{question}
\gkti{9}双曲线 $2mx^2-my^2=2$ 的一条准线是 $y=1$，则 $m=$\gkblank{}。

\end{question}
\begin{question}
\gkti{10}平面上，四条平行直线与另外五条平行直线互相垂直，则它的矩形共有\gkblank{}个。（结果用数值表示）

\end{question}
\gkpart{二、选择题}

\begin{question}
\gkti{11}圆的半径是 $1$，圆心的极坐标是 $(1,0)$，则这个圆的极坐标方程是\mbox{（　　）}。

（A）$\rho=\cos\theta$　　（B）$\rho=\sin\theta$　　（C）$\rho=2\cos\theta$　　（D）$\rho=2\sin\theta$

\end{question}
\begin{question}
\gkti{12}函数 $f(x)$ 和 $g(x)$ 的定义域均为 $\mathbb R$，“$f(x),g(x)$ 都是奇函数”是“$f(x)$ 与 $g(x)$ 的积是偶函数”的\mbox{（　　）}。

（A）必要条件但非充分条件　　（B）充分条件但非必要条件

（C）充分必要条件　　（D）非充分条件也非必要条件

\end{question}
\begin{question}
\gkti{13}设点 $P$ 在有向线段 $\overrightarrow{AB}$ 的延长线上，$P$ 分 $\overrightarrow{AB}$ 所成的比为 $\lambda$，则\mbox{（　　）}。

（A）$\lambda<-1$　　（B）$-1<\lambda<0$　　（C）$0<\lambda<1$　　（D）$\lambda>1$

\end{question}
\begin{question}
\gkti{14}设 $2^a=3$，$2^b=6$，$2^c=12$，则数列 $a,b,c$\mbox{（　　）}。

（A）是等差数列但不是等比数列　　（B）是等比数列但不是等差数列

（C）既是等差数列又是等比数列　　（D）既不是等差数列又不是等比数列

\end{question}
\begin{question}
\gkti{15}设 $\alpha$ 角属于第二象限，且 $\left|\cos\dfrac\alpha2\right|=-\cos\dfrac\alpha2$，则 $\dfrac\alpha2$ 角属于\mbox{（　　）}。

（A）第一象限　　（B）第二象限　　（C）第三象限　　（D）第四象限

\end{question}
\begin{question}
\gkti{16}设过长方体同一个顶点的三个面的对角线长分别是 $a,b,c$，那么这个长方体的对角线长是\mbox{（　　）}。

（A）$\sqrt{a^2+b^2+c^2}$　　（B）$\sqrt{\dfrac{a^2+b^2+c^2}{2}}$

（C）$\sqrt{\dfrac{a^2+b^2+c^2}{3}}$　　（D）$\dfrac{\sqrt{a^2+b^2+c^2}}{2}$

\end{question}
\begin{question}
\gkti{17}函数 $f(x)=a\tan\dfrac xa$ 的最小正周期是\mbox{（　　）}。

（A）$\pi a$　　（B）$\pi|a|$　　（C）$\dfrac\pi a$　　（D）$\dfrac\pi{|a|}$

\end{question}
\begin{question}
\gkti{18}已知 $1<x<d$，令 $a=(\log_d x)^2$，$b=\log_d(x^2)$，$b=\log_d(\log_d x)$，则\mbox{（　　）}。

（A）$a<b<c$　　（B）$a<c<b$　　（C）$c<b<a$　　（D）$c<a<b$

\end{question}
\begin{question}
\gkti{19}设 $a,b$ 是两条异面直线，那么下列四个命题中的假命题是\mbox{（　　）}。

（A）经过直线 $a$ 有且只有一个平面平行于直线 $b$

（B）经过直线 $a$ 有且只有一个平面垂直于直线 $b$

（C）存在分别经过直线 $a$ 和 $b$ 的两个互相平行的平面

（D）存在分别经过直线 $a$ 和 $b$ 的两个互相垂直的平面

\end{question}
\begin{question}
\gkti{20}下列四个函数中，在定义域内不具有单调性的函数是\mbox{（　　）}。

（A）$y=\cot(\arccos x)$　　（B）$y=\tan(\arcsin x)$

（C）$y=\sin(\arctan x)$　　（D）$y=\cos(\arctan x)$

\end{question}
\gkpart{三、解答题}

\begin{question}
\gkti{21}已知 $\log_5(x^2+2x-2)=0$，$2\log_5(x+2)-\log_5 y+\dfrac12=0$，求 $y$ 的值。

\end{question}
\begin{question}
\gkti{22}求方程 $\sqrt{5\cos x+\cos 2x}+\sin x=0$ 在 $[0,2\pi)$ 上的解。

\end{question}
\begin{question}
\gkti{23}已知点 $P$ 直线 $x=2$ 上移动，直线 $l$ 通过原点且与 $OP$ 垂直；通过点 $A(1,0)$ 及点 $P$ 的直线 $m$ 和直线 $l$ 交于点 $Q$。求点 $Q$ 的轨迹方程，并指出该轨迹的名称和它的焦点坐标。

\end{question}
\begin{question}
\gkti{24}已知直线 $l:x-ny=0$（$n\in\mathbb N$）；圆 $M:(x+1)^2+(y+1)^2=1$；抛物线 $\Phi:y=(x-1)^2$。又 $L$ 与 $M$ 交于点 $A,B$；$L$ 与 $\Phi$ 交于点 $C,D$。求
\nopagebreak[4]
\[
\lim_{n\to\infty}\frac{|AB|^2}{|CD|^2}.
\]

\end{question}
\begin{question}
\gkti{25}关于实数 $x$ 的不等式
\nopagebreak[4]
\[
\left|x-\frac{(a+1)^2}{2}\right|\leq\frac{(a-1)^2}{2}
\]
\Needspace{5\baselineskip}
与
\nopagebreak[4]
\[
x^2-3(a+1)x+2(3a+1)\leq0
\]
的解集依次记为 $A$ 与 $B$。求使 $A\subseteq B$ 的 $a$ 的取值范围（其中 $a\in\mathbb R$）。

\Needspace{0.4\textheight}
\end{question}
\begin{question}
\gkti{26}如图，平面 $\alpha,\beta$ 相交于直线 $MN$，点 $A$ 在平面 $\alpha$ 上，点 $B$ 在平面 $\beta$ 上，点 $C$ 在直线 $MN$ 上，$\angle ACM=\angle BCN=45^\circ$，$A$—$MN$—$B$ 是 $60^\circ$ 的二面角，$AC=1$。求：

（1） 点 $A$ 到平面 $\beta$ 的距离；

（2） 二面角 $A$—$BC$—$M$ 的大小（用反三角函数表示）。
\nopagebreak[4]
\begin{center}% 1990 上海卷第 26 题：两平面及点线拓扑按原卷重绘，不表达精确比例。
\begin{tikzpicture}[x=0.78cm,y=0.74cm,line cap=round,line join=round]
  \coordinate (M) at (0,0);
  \coordinate (N) at (6,0);
  \coordinate (C) at (2.55,0);
  \coordinate (A) at (3.15,1.05);
  \coordinate (B) at (4.45,-0.9);
  \coordinate (AlphaLeft) at (1.25,1.55);
  \coordinate (AlphaRight) at (7.25,1.55);
  \coordinate (BetaLeft) at (1.2,-1.45);
  \coordinate (BetaRight) at (7.2,-1.45);

  % 两平面共用 MN；远侧的 alpha 边按原图用虚线表示。
  \draw[thin] (M)--(N)--(AlphaRight);
  \draw[thin,dashed] (AlphaRight)--(AlphaLeft);
  \draw[thin] (AlphaLeft)--(M);
  \draw[thin] (N)--(BetaRight)--(BetaLeft)--(M);

  % 原图中 A、B、C 的连线拓扑；透视坐标仅为排布，不代表测量值。
  \draw[thick] (A)--(C) (C)--(B) (A)--(B);

  \node[below left=2pt] at (M) {$M$};
  \node[below right=2pt] at (N) {$N$};
  \node[below=4pt] at (C) {$C$};
  \node[above left=1pt] at (A) {$A$};
  \node[below right=1pt] at (B) {$B$};
  \node[above right=3pt] at (AlphaRight) {$\alpha$};
  \node[below=5pt] at (BetaLeft) {$\beta$};
  \node[font=\scriptsize,align=center] at (3.6,-2.15)
    {按原卷点线关系重绘；图形不按比例。};
\end{tikzpicture}
\end{center}
\nopagebreak[4]
\end{question}
\begin{question}
\gkti{27}复平面上点 $A,B$ 对应的复数分别为 $z_1=2,z_2=-3$，点 $P$ 对应的复数为 $z$，$\dfrac{z-z_1}{z-z_2}$ 的辐角主值为 $\varphi$。当点 $P$ 在以原点为圆心、$1$ 为半径的上半圆周（不包括两个端点）上运动时，求 $\varphi$ 的最小值。

\end{question}